% GENERATED FILE. Edit canonical lessons, then run npm run generate:books.
% canonical-source-hash: fnv1a64:e6600d7a15b29d11
% canonical-lessons: MR-C252-maja-karne, MR-C252-ragavne, MR-C252-kalji-karne, MR-C252-dhumrapan-karne, MR-C252-sobat-jane

\chapter{To have fun, To get angry, To worry, To smoke, To go with}
\label{ch:mr-to-have-fun-to-get-angry-to-worry-to-smoke-to-go-with}
\hlchaptermodality{252}

\begin{chapteropening}
\textbf{By the end of this chapter:} \emph{I can say to have fun, to get angry, to worry, to smoke and to go with.}

Complete the last lesson of chapter 252: I can say to have fun, to get angry, to worry, to smoke and to go with.
\end{chapteropening}

\section[\texorpdfstring{\mr{मजा} \mr{करणे} (majā karṇe)}{मजा करणे (majā karṇe)}]{\mr{मजा} \mr{करणे} (majā karṇe) — to have fun}

\label{lesson:MR-C252-maja-karne}



\begin{quote}
Before the new one: say the Marathi for to move, then the Marathi for to heat.
\end{quote}

\subsection*{You'll want to know: \mr{मजा} \mr{करणे}}
\textbf{\mr{मजा} \mr{करणे}} — \emph{majā karṇe} — \textquotedblleft{}to have fun\textquotedblright{}.

\begin{grammarlens}[title={What you've built}]
One more word for this chapter.
\end{grammarlens}

\subsection*{Your turn}
\begin{itemize}
\raggedright
  \item \emph{Say it:} \emph{majā karṇe}
  \item \emph{Say it:} \emph{majā karṇe}, once more
  \item \emph{Say it:} say \emph{majā karṇe}
  \item \emph{Recall:} say the Marathi for to move, then the Marathi for to heat, then say \emph{majā karṇe} again
\end{itemize}

\begin{tcolorbox}[breakable,colback=teal!4,colframe=teal!35!black,title={Before you move on}]
What does \mr{मजा} \mr{करणे} mean? (\textquotedblleft{}To have fun\textquotedblright{}.) Say it once more.
\end{tcolorbox}
\section[\texorpdfstring{\mr{रागावणे} (rāgāvṇe)}{रागावणे (rāgāvṇe)}]{\mr{रागावणे} (rāgāvṇe) — to get angry}

\label{lesson:MR-C252-ragavne}



\begin{quote}
Before the new one: say the Marathi for to heat, then the Marathi for to have fun.
\end{quote}

\subsection*{You'll want to know: \mr{रागावणे}}
\textbf{\mr{रागावणे}} — \emph{rāgāvṇe} — \textquotedblleft{}to get angry\textquotedblright{}.

\begin{grammarlens}[title={What you've built}]
One more word for this chapter.
\end{grammarlens}

\subsection*{Your turn}
\begin{itemize}
\raggedright
  \item \emph{Say it:} \emph{rāgāvṇe}
  \item \emph{Say it:} \emph{rāgāvṇe}, once more
  \item \emph{Say it:} say \emph{rāgāvṇe}
  \item \emph{Recall:} say the Marathi for to heat, then the Marathi for to have fun, then say \emph{rāgāvṇe} again
\end{itemize}

\begin{tcolorbox}[breakable,colback=teal!4,colframe=teal!35!black,title={Before you move on}]
What does \mr{रागावणे} mean? (\textquotedblleft{}To get angry\textquotedblright{}.) Say it once more.
\end{tcolorbox}
\section[\texorpdfstring{\mr{काळजी} \mr{करणे} (kāḷjī karṇe)}{काळजी करणे (kāḷjī karṇe)}]{\mr{काळजी} \mr{करणे} (kāḷjī karṇe) — to worry}

\label{lesson:MR-C252-kalji-karne}



\begin{quote}
Before the new one: say the Marathi for to have fun, then the Marathi for to get angry.
\end{quote}

\subsection*{You'll want to know: \mr{काळजी} \mr{करणे}}
\textbf{\mr{काळजी} \mr{करणे}} — \emph{kāḷjī karṇe} — \textquotedblleft{}to worry\textquotedblright{}.

\begin{grammarlens}[title={What you've built}]
One more word for this chapter.
\end{grammarlens}

\subsection*{Your turn}
\begin{itemize}
\raggedright
  \item \emph{Say it:} \emph{kāḷjī karṇe}
  \item \emph{Say it:} \emph{kāḷjī karṇe}, once more
  \item \emph{Say it:} say \emph{kāḷjī karṇe}
  \item \emph{Recall:} say the Marathi for to have fun, then the Marathi for to get angry, then say \emph{kāḷjī karṇe} again
\end{itemize}

\begin{tcolorbox}[breakable,colback=teal!4,colframe=teal!35!black,title={Before you move on}]
What does \mr{काळजी} \mr{करणे} mean? (\textquotedblleft{}To worry\textquotedblright{}.) Say it once more.
\end{tcolorbox}
\section[\texorpdfstring{\mr{धूम्रपान} \mr{करणे} (dhūmrapān karṇe)}{धूम्रपान करणे (dhūmrapān karṇe)}]{\mr{धूम्रपान} \mr{करणे} (dhūmrapān karṇe) — to smoke}

\label{lesson:MR-C252-dhumrapan-karne}



\begin{quote}
Before the new one: say the Marathi for to get angry, then the Marathi for to worry.
\end{quote}

\subsection*{You'll want to know: \mr{धूम्रपान} \mr{करणे}}
\textbf{\mr{धूम्रपान} \mr{करणे}} — \emph{dhūmrapān karṇe} — \textquotedblleft{}to smoke\textquotedblright{}.

\begin{grammarlens}[title={What you've built}]
One more word for this chapter.
\end{grammarlens}

\subsection*{Your turn}
\begin{itemize}
\raggedright
  \item \emph{Say it:} \emph{dhūmrapān karṇe}
  \item \emph{Say it:} \emph{dhūmrapān karṇe}, once more
  \item \emph{Say it:} say \emph{dhūmrapān karṇe}
  \item \emph{Recall:} say the Marathi for to get angry, then the Marathi for to worry, then say \emph{dhūmrapān karṇe} again
\end{itemize}

\begin{tcolorbox}[breakable,colback=teal!4,colframe=teal!35!black,title={Before you move on}]
What does \mr{धूम्रपान} \mr{करणे} mean? (\textquotedblleft{}To smoke\textquotedblright{}.) Say it once more.
\end{tcolorbox}
\section[\texorpdfstring{\mr{सोबत} \mr{जाणे} (sobat jāṇe)}{सोबत जाणे (sobat jāṇe)}]{\mr{सोबत} \mr{जाणे} (sobat jāṇe) — to go with}

\label{lesson:MR-C252-sobat-jane}



\begin{quote}
Before the new one: say the Marathi for to worry, then the Marathi for to smoke.
\end{quote}

\subsection*{You'll want to know: \mr{सोबत} \mr{जाणे}}
\textbf{\mr{सोबत} \mr{जाणे}} — \emph{sobat jāṇe} — \textquotedblleft{}to go with\textquotedblright{}.

\begin{grammarlens}[title={What you've built}]
One more word for this chapter.
\end{grammarlens}

\subsection*{Your turn}
\begin{itemize}
\raggedright
  \item \emph{Say it:} \emph{sobat jāṇe}
  \item \emph{Say it:} \emph{sobat jāṇe}, once more
  \item \emph{Say it:} say \emph{sobat jāṇe}
  \item \emph{Recall:} say the Marathi for to worry, then the Marathi for to smoke, then say \emph{sobat jāṇe} again
\end{itemize}

\begin{tcolorbox}[breakable,colback=teal!4,colframe=teal!35!black,title={Before you move on}]
What does \mr{सोबत} \mr{जाणे} mean? (\textquotedblleft{}To go with\textquotedblright{}.) Say it once more.
\end{tcolorbox}
